\section{公司的设立}
\begin{enumerate}
  \item 设立。发起、募集。
        \begin{enumerate}
          \item 发起人/创始股东（股东不一定是发起人）。起草、签署章程；出钱；履行设立职责。
                \begin{enumerate}
                  \item 有限 1 - 50 人；股份 2 - 200 + 半数在国内。\bs{视为民事合伙关系}。
                  \item 人数总结。普通合伙 2+、有限合伙 2 - 50、个人独资 1、个体工商 1 - 8。
                  \item \bs{侵权处理}。
                        \begin{enumerate}
                          \item \zx{以个人名义}做事。\zx{成立后择一行权，但不得主张连带责任}。
                          \item 以\zx{设立中公司名义} \yjt 找公司。\zd{为个人利益 + 相对人知情 \yjt 找个人，公司免责}。
                          \item \zd{失败。费用、债务大家对外连带，对内约定 \yjt 出资比例 \yjt 平均}。
                          \item 费用。\zd{发起人过错的发起人承担}，其他合理费用公司承担。
                          \item \sjt \bs{内外有别，看谁过错}。
                        \end{enumerate}
                \end{enumerate}
        \end{enumerate}
  \item 成立。签发营业执照。
  \item 名称。行政区划、字号、行业 \hz 经营特点、组织形式。
  \item 章程。必备要件。登记对抗，但不是生效要件。
        \begin{enumerate}
          \item 强制事项、任意事项。
          \item 只对公司、股东、董监高具有约束力。
        \end{enumerate}
  \item \bs{股东出资制度}。必考。理论上应坚持对公转账，打到公司账户。实践中，执行很难，乱象很多。
        \begin{enumerate}
          \item 出资形式。
                \begin{enumerate}
                  \item 货币。来者不拒。非法不能直接划走，只能通过拍卖 \hz 变卖处置股权。\al{甲 贪污 50W 出资}。
                  \item \zx{非货币}。\zd{可估价 + 要转移所有权（土地除外）}。\al{厂房、设备；甲以名下房产 20 年\xhx{使用权}出资 \cuo}。
                        \begin{enumerate}
                          \item 债权。\bs{本质是债权让与}。\al{乙以其对丙估值 300W 的债权出资，后债权到期变为 30W}。
                          \item 房屋、土地使用权、知识产权。应\bs{及时交付 + 登记}。\zx{交付视为移转所有权}，股东开始享权。
                          \item 抵押地/划拨地 \yjt 合理期限\zd{解除权利负担}/办理土地变更手续。
                          \item 股权。要走法定手续。
                        \end{enumerate}
                  \item \sjt \bs{无权处分，适用善意取得}。
                        \al{甲乙兄弟，房继承给乙，甲伪造遗嘱，并以房子出资和丙（知情）、丁创办公司，担任总经理。\xhx{其他人知道、担任高管职位 \yjt 公司视为知情}}。
                  \item 禁止出资的形式。劳务、信用、商誉、特许经营、设定担保的财产。不好估价 \hz 转移不了所有权。
                \end{enumerate}
          \item 出资方式。无最低注册资本、出资期限要求。
                \begin{enumerate}
                  \item 认缴。投票按认缴比例来。实收资本 \yjt 代表当年情况。总资产 \yjt 代表现在情况。
                  \item 实缴。金融类公司。
                \end{enumerate}
          \item \bs{（一次性）出资责任}。
                \begin{enumerate}
                  \item 对内。
                        \begin{enumerate}
                          \item \zx{出资违约}（仅限货币财产出资）。未按期足额，\zd{补足 + 违约责任 + 发起人连带}。
                          \item \zx{出资不实}。\zd{补足 + 发起人连带}。非股东原因导致股票/债权价值下跌的不算出资不实。
                          \item \zx{抽逃}。\zd{虚假财务会计报表、虚增利润分配、虚构债权债务转出、关联交易} \yjt 实质损害。
                                \begin{itemize}
                                  \item \zd{抽逃者返还本息}，\zx{协助者连带}。
                                  \item 超出部分认定为侵害公司的财产权进行处理。
                                \end{itemize}
                        \end{enumerate}
                  \item 对外。债权人。
                        \begin{enumerate}
                          \item 行使代位权。
                          \item \bs{对不能清偿的债务，瑕疵股东补充责任、发起人连带责任}、高管（增资时）。
                                \al{债权人可要甲\xhx{补足出资} \cuo{债权人要的是钱，出资的钱不到他手上；出资是\xhx{内部的事}，管不着}}。
                        \end{enumerate}
                \end{enumerate}
          \item 出资期限利益。原则上不能被\zd{加速到期}。4 个例外。不想还钱、还不起钱。
                \begin{enumerate}
                  \item 公司解散、破产申请被受理。
                  \item 被执行 + 没有财产可供执行 + 该破不破。
                  \item 恶意逃债行为。\al{债务产生后，股东会决议 \hz 其他方式延长股东出资期限}。
                \end{enumerate}
          \item \zx{出资瑕疵。不能直接认定无效，合理期限 + 补正措施}。权利限制。\al{剥夺 \cuo}。
                \begin{enumerate}
                  \item 新、剩、利。合理限制新股优先认购、剩余财产分配、利润分配。
                  \item 没出 \yjt 催告 + 股东会议 \yjt 解除股东资格 \yjt 减资 \hz 补。\al{\xhx{直接}开会表决解除 \cuo}。
                \end{enumerate}
        \end{enumerate}
\end{enumerate}
